%------------------------------------------------------------------------------
% CHAPTER 2: OVERALL DESCRIPTION
%------------------------------------------------------------------------------
\chapter{Overall Description}
\label{ch:overall}

\section{Product Perspective}
\label{sec:perspective}

Best Bike Paths (\texttt{BBP}) is a standalone software meant to help cyclists facilitate the creation and management of a collective bike paths inventory. The software operates as a distributed system (like a downloadable smartphone application) that integrates both user-generated content and automated sensor data retrieved from users' mobile device/devices.

Specifically \texttt{BBP} interacts with mobile hardware interfaces, leveraging \texttt{GPS} to reconstruct trip trajectories and inertial sensors (gyroscopes and accelerometers) to infer road surfaces conditions and detect anomalies such as potholes. Furthermore, the system interfaces with external meteorological services to enrich trip data with real-time variables like temperature, wind-speed and raining conditions.

The system provides accurate and up-to-date details regarding paths thanks to a data aggregation engine capable of merging conflicting information from multiple sources based on freshness and confirmation from users.

\texttt{BBP} functions as an independent entity for the user but it relies heavily on this ecosystem of mobile sensors and external \texttt{APIs} (weather forecast, map display and so on).

\subsection{Scenarios}
\label{sec:scenarios}

\begin{enumerate}
    \item \textbf{User Signup}: The guest user Alice accesses the sign-up interface to create a profile. Upon successful validation of the provided registration details, the system grants the new user the "Registered" status, allowing him to access new functions reserved to Registered Users.

    \item \textbf{User login}: The registered user Bob inputs his credentials in \texttt{BBP} which validates them and loads the user's personal session.

    \item \textbf{Automated trip recording}: The registered user Carol initiates a recording session. Once Carol begins riding, the \texttt{BBP} system automatically detects the cycling activity and begins logging GPS coordinates, sensor data, and timestamps. Recording continues until Carol pauses cycling or stops the session.
    
    \item \textbf{Anomaly validation}: Following an automated trip, the user Carol has 72 hours to review potential obstacles detected through \texttt{BBP} sensor data analysis and confirm or reject identified potholes, preventing false positives from entering the database.

    \item \textbf{Trip Publishing}: The registered user Carol accesses her personal history to select a trip recorded and validated in the past to share the collected data with the community. \texttt{BBP} ensures data anonymization before publishing it

    \item \textbf{Private Trip Reviews}: The registered user Carol accesses one of her private trips to visualize it on the map or view recording statistics, such as average speed, trip duration and weather conditions at the time of recording.

    \item \textbf{Manual Path Insertion}: The registered user David manually inserts information about a specific bike path by entering individual street names and their status (e.g., "optimal", "requires maintenance") without requiring physical cycling or automated recording.

    \item \textbf{Path searching}: The user Erwin, whether registered or not, enters an origin and a destination. She can do it by either inserting a street name or by selecting a point on the map. The \texttt{BBP} computes and visualizes the available paths connecting the two chosen endpoints. When multiple routes exist for a given query, the user Erwin will view them ordered by their "Path score".

    \item \textbf{Path Score Calculation}: The registered user Frank publishes a trip reporting a bike path as \texttt{OPTIMAL}. The \texttt{BBP} system merges Erwin's trip with existing data from other users who previously reported the same path as \texttt{REQUIRES\_MAINTENANCE}. Based on data freshness and consensus, \texttt{BBP} recalculates the path status and updates the path score considering both the new aggregated status and the path's effectiveness (average trip time).

    \item \textbf{Path Details Viewing}: The user Frank, whether registered or not, selects a bike path from search results to view detailed information including current status, trajectory on the map and current weather conditions along the path.
    
\end{enumerate}


\subsection{Class Diagram}
\label{sec:class_diagram}

The \texttt{UML} Class diagram below represents a high level model of the mobile application described in this document. There are two kinds of users: \texttt{RegisteredUser} and its superclass \texttt{User}. \texttt{Trip} class represents a bike ride of a \texttt{RegisteredUser}, it can be an \texttt{AutomaticTrip} or a \texttt{ManualTrip}. \\\\
Whenever a \texttt{Trip} is published by its creator, the system anonymizes and merges it with other \texttt{Trip}(s) from different \texttt{RegisteredUser}(s), contributing to the public \texttt{Path}, which can be searched, found and seen on map by \texttt{User}(s). 
\\\\
The \texttt{ManualTrip} class has a list of streets, which are inserted by the \texttt{RegisteredUser}. They are suddenly translated into \texttt{GPSCoordinates} by using an external \texttt{MapService}, also used for \texttt{Street} lookup and \texttt{Path} visualization. \\
On the other hand, \texttt{AutomaticTrip} leverages the \texttt{MobileSensors} for directly obtaining \texttt{GPSCoordinates} and \texttt{SensorReadings} which are analyzed to find the \texttt{Obstacle}(s). Any \texttt{Obstacle} must be confirmed by the \texttt{RegisteredUser} recording the \texttt{Trip} before it can be published. \texttt{Obstacle}s can be aggregated into \texttt{Path} during the merge of multiple \texttt{Trip}s
\\\\
\texttt{Statistics} depending on speed, such as average speed and trip time, are only recorded for \texttt{AutomaticTrip}. \texttt{Weather} forecasts are provided by an external \texttt{WeatherService} and are used within \texttt{Statistics} since they can influence performance. Also, any \texttt{User} consults the \texttt{WeatherService} when it views a public \texttt{Path}.

\begin{figure}[H]
    \centering
    \includesvg[width=\textwidth]{diagrams/class_diagram.svg}
    \caption{Class Diagram}
    \label{fig:class_diagram}
\end{figure}

\subsection{State Diagrams}
\label{sec:state_diagrams}

\subsubsection{Path}

This state diagram shows how the status of a \texttt{Path} evolves as \texttt{BBP} receives already-anonymized publishable information from different Registered Users. When a first \texttt{Trip} is published and has no corresponding \texttt{Path}, an instance is created, corresponding to the \texttt{CREATED} state. Once the first published \texttt{Trip} is merged, during the \texttt{MERGE\_DATA} stage, the \texttt{Path} enters one of the status states (e.g., \texttt{OPTIMAL} or \texttt{REQUIRES\_MAINTENANCE}).\\
When one or more \texttt{Trip}(s) corresponding to the same path are published, their information is merged into the corresponding \texttt{Path}. The \texttt{Path} status reflects the status of the most recently published \texttt{Trip}. If more than two \texttt{Trip}(s) referring to the same \texttt{Path} are published in the same timeframe window, the Path status is determined by majority voting among their reported statuses. If \texttt{BBP} couldn't find a strict majority, it keeps the previous status. Also, during the merge process, \texttt{BBP} aggregates obstacles from incoming trips into the path.

\texttt{Trip} then moves to the newly computed path status corresponding state \\
When all reliable information indicates that the path can no longer be used, it moves to the \texttt{PERMANENTLY\_CLOSED} final state. As the new status, also the new path score is recomputed during the \texttt{MERGE\_DATA} state.

\begin{figure}[H]
    \centering
    \includesvg[width=0.8\textwidth]{diagrams/path_state_diagram.svg}
    \caption{Path State Diagram}
    \label{fig:path_state}
\end{figure}

\subsubsection{Manual Trip}

This state diagram describes the lifecycle of a \texttt{ManualTrip}. When a Registered User clicks on the manual \texttt{Trip} recording button, the trip enters the \texttt{INSERTING\_STREET} state, during which the user specifies the list of streets and assigns a status to each segment one at a time. The user can iteratively add streets until satisfied. Once the user submits the form, \texttt{BBP} automatically computes intersections between streets and computes the total distance, saving the \texttt{Trip} in the \texttt{SAVED\_PRIVATE} state. Unlike \texttt{AutomaticTrip}, \texttt{ManualTrip} does not involve obstacle detection, so the \texttt{Trip} automatically transitionsto \texttt{READY\_TO\_PUBLISH} without requiring validation. When the user decides to share the collected information with the community, the trip moves to the \texttt{PUBLISHED} state, triggering the \texttt{Path.MERGE\_DATA} state and contributing to its corresponding public path. From now on, at any time, the user might delete the \texttt{Trip}, moving it to the \texttt{DELETED} final state.

\begin{figure}[H]
    \centering
    \includesvg[width=0.8\textwidth]{diagrams/manual_trip_state_diagram.svg}
    \caption{Manual Trip State Diagram}
    \label{fig:manual_trip_state}
\end{figure}

\subsubsection{Automatic Trip}

The lifecycle of an \texttt{AutomaticTrip} begins in \texttt{IDLE} state. Once the user clicks on the "Start automatic recording" button, it moves to \texttt{DETECT\_BIKING}. In this state, the \texttt{BBP} monitors movement: if the user speed exceeds the biking speed threshold, it enters the \texttt{RECORDING} state, during which GPS coordinates and sensor data are collected; if it drops below, it reverts to \texttt{DETECT\_BIKING}. After the user stops the recording, the \texttt{Trip} undergoes \texttt{ANALYSIS}, where \texttt{BBP} analyzes sensor data to identify obstacles, computes statistics and queries the external \texttt{Weather Service} to retrieve meteorogical data, afterwards moving to the \texttt{SAVED\_PRIVATE} state. \\ 
Then, If no obstacles are found, it proceeds to the \texttt{READY\_TO\_PUBLISH} state. Otherwise, it enters the \texttt{AWAITING\_VALIDATION} stage. From here, the user has 72 hours to validate the data to reach \texttt{READY\_TO\_PUBLISH}; otherwise, the \texttt{Trip} becomes \texttt{NOT\_PUBLISHABLE}. Finally, the user can transition the \texttt{Trip} to \texttt{PUBLISHED} or, from any of the previous states, to \texttt{DELETED}, or simply keep it in their private history to end the process.

\begin{figure}[H]
    \centering
    \includesvg[width=\textwidth]{diagrams/automatic_trip_state_diagram.svg}
    \caption{Automatic Trip State Diagram}
    \label{fig:automatic_trip_state}
\end{figure}


\section{Product Functions}
\label{sec:functions}

\begin{itemize}
    \item \textbf{User Management}: Allows unregistered users to sign up and registered users to log in and access rights to restricted features like trip recording and publishing.

    \item \textbf{Automated Trip Recording}: Automatically detects when a registered user is cycling based on mobile sensors' data. It records the trajectory via GPS and continuously monitors device sensors (accelerometer, gyroscope) to identify road anomalies and significant movements.

    \item \textbf{Manual Activity Tracking}: Enables registered users to manually insert bike path information by specifying street names and assigning them a status without the need of physical cycling and real-time tracking.

    \item \textbf{Anomaly detection and user validation}: Processes sensor data to infer the presence of obstacles (e.g., potholes). It presents these potential anomalies to the user for post-trip-recording validation to avoid false positives.

    \item \textbf{Data Enrichment}: Integrates with external weather services to enrich trip data with weather conditions, temperature and wind speed. This involves providing historical weather conditions to past trips and real-time weather conditions along the path a user is currently visualizing.

    \item \textbf{Data Merging}: Implements a merge process that fuses information about the same path from different users. It resolves conflicts and determines the current path status based on freshness and number of confirming data obtained.

    \item \textbf{Path Score Calculation}: Automatically computes a quality metric for bike paths based on aggregated status and route effectiveness (average trip time), recalculating whenever new trip data is merged.

    \item \textbf{Path Search}: Enables users to query routes between two points and view them on a map, ranked by their path scores to facilitate informed route selection.

    \item \textbf{Statistical Reporting}: Generates performance metrics for registered users, offering the possibility to review past trips with data such as average speed and total trip time.
\end{itemize}

\section{User Characteristics}
\label{sec:users}

In the BBP platform, there are two kind of user: the generic one \texttt{User} and the \texttt{Registered User}

\subsubsection{ User}

A User can explore bike paths without creating an account.

He can search for paths between any chosen origin and destination and view the corresponding trajectory directly on the map, selecting the safest (path status) or most efficient one (effectiveness = average trip time). In fact, if multiple paths exist, \texttt{BBP} shows them to the user ranked by a score based on the path status and the effectiveness of the path in reaching the destination from the given origin. For each proposed path the User can inspect its public characteristics, such as total distance, estimated trip time, and the status validated by the community.

Users benefit from always-updated information, since \texttt{BBP} merges multiple contributions provided by different registered users to determine the most reliable and fresh status of each path.

Generic users have read-only access to all public information collected by the system and cannot record trips, validate anomalies, or store personal path histories. Their interaction with \texttt{BBP} mainly focuses on quickly retrieving reliable and up-to-date information about the best available routes for their ride.

\subsubsection{Registered User}

A Registered User can access all the features available to the generic User, but he can also access more advanced functionalities and actively contribute to the \texttt{BBP} system. In particular, he can record his trips and store them on the platform, keeping track of his cycling activity and contributing to the public bike paths inventory.

He can record his trips in two different ways:
\begin{itemize}
    \item \textbf{Manual mode}: by inserting manually the names of the streets belonging to a bike path and assigning a status to each street (such as \texttt{optimal}, \texttt{medium}, \texttt{sufficient} and \texttt{requires maintenance}).
    \item \textbf{Automatic mode}: letting \texttt{BBP} acquire data from their mobile device sensors, detecting when he starts cycling, collecting GPS data to reconstruct the followed trajectory, and using accelerometer and gyroscope data to identify potential anomalies such as potholes. After the trip, the Registered User must confirm or reject the automatically detected anomalies within 72 h, before they become publishable, otherwise they will never be published.
\end{itemize}

A Registered User can view his personal cycling history, accessing stored trips enriched with statistics such as average speed, trip time and weather conditions at the time of recording.

When desired, he can publish one of his validated bike paths, including information such as total distance, estimated trip time and the status, triggering the merge process into a public path and making its information available to all other users.

\section{Assumptions, Dependencies and Constraints}
\label{sec:assumptions}

\subsection{Domain Assumptions}

\begin{itemize}
    \item \textbf{[D.1.1] Sensor Reliability}: It is assumed that GPS, accelerometer and gyroscope sensors provided by the mobile hardware are sufficiently precise to distinguish biking vibrations from road anomalies effectively.

    \item \textbf{[D.1.2] Biking Speed Distinctiveness}: It is assumed that the speed and the acceleration of a user cycling are mathematically distinguishable from those of walking, running or driving a car, enabling the Automated Mode to function correctly.

    \item \textbf{[D.1.3] User Trustworthiness}: Given the community-based nature of the data, it is assumed that Registered Users will generally provide truthful validations regarding obstacles and path status.

    \item \textbf{[D.1.4] Obstacle Detection}: If different Automatic Trips corresponding to the same Path detect the same obstacle, we assume that the obstacle has the same coordinates in both trips.
    
\end{itemize}

\subsection{Dependencies}

\begin{itemize}
    \item \textbf{[D.2.1] External Meteorological Service}: The system relies on third-party weather APIs to retrieve environment conditions. The availability and accuracy of "weather-enriched" statistics depend solely on these external provider/s.

    \item \textbf{[D.2.2] Mapping and Geocoding Service}: The visualization of paths and the translation of GPS coordinates to street names (Reverse Geocoding) depend on external map providers (e.g., Google Maps, OpenStreetMap).

    \item \textbf{[D.2.3] Mobile Hardware Ecosystem}: The application is strictly dependent on the user's device having functional hardware sensors and an OS that grants the necessary permissions for background tracking.
\end{itemize}

\subsection{Constraints}

\begin{itemize}
    \item \textbf{[C.1] Network Connectivity}: A continuous network connection is required to access maps, fetch weather data and synchronize shared path information. Offline functionality is constrained to local buffering of the current trip recording.

    \item \textbf{[C.2] Power Consumption}: The continuous polling of GPS and motion sensors is resource-intensive. The software implementation must be optimized to prevent excessive battery drainage, which would hinder the user's ability to record long trips.

    \item \textbf{[C.3] Regulatory Compliance}: The storage of geolocation data and user habits is subject to strict privacy regulations. The system must enforce explicit consent mechanisms for data tracking and sharing.
\end{itemize}